\section{Perspectiva práctica: una lista de verificación de tres capas para un Coding Agent mantenible}

\subsection{Lo que más se pasa por alto es la gobernanza de la sesión y el contexto}
La primera mitad de la semana 2 usó un Coding Agent mínimo para mostrar la estructura central del tool loop, y la segunda mitad introdujo MCP. Al ver las dos partes juntas, aparece un problema de ingeniería más práctico: \textbf{en un Coding Agent que funcione de manera estable a largo plazo, lo decisivo no es solo si sabe invocar herramientas, sino cómo administra el contexto, el estado y la recuperación tras fallas.}

\begin{knowledgebox}{Lista de verificación de tres capas}
\begin{enumerate}
    \item \textbf{Capa del modelo}: si el prompt del sistema explica con claridad la tarea, los límites y los principios de uso de las herramientas;
    \item \textbf{Capa de las herramientas}: si el schema de las herramientas es estable y si los resultados que regresan son adecuados para el siguiente paso de razonamiento;
    \item \textbf{Capa de la sesión}: si el contexto tiende a inflarse, si existen reminders, y cómo se recupera tras una falla.
\end{enumerate}
\end{knowledgebox}

Esto también explica por qué el curso insiste en la precarga del context, los system reminders y la derivación hacia subagents. Estas cosas que parecen «trucos de producto» en realidad compensan la fragilidad que tienen los modelos actuales en tareas de largo plazo. No basta con que el Agent sepa invocar herramientas: también debe mantener la coherencia con el objetivo durante un trabajo largo, con muchos pasos y muchos archivos.

\subsection{Resumen de la sección}
Si lo aprendido en la semana 2 se lleva de verdad a la práctica de ingeniería, un Coding Agent confiable debe superar al mismo tiempo al menos tres pruebas: si el modelo sabe qué hacer, si las herramientas son adecuadas para que el modelo las invoque, y si la sesión cuenta con suficiente capacidad de gobernanza del contexto. Ninguna de las tres puede faltar.

\subsection{Cuándo no conviene reducir el problema a «agregar otro MCP Server»}
La semana 2 deja otro juicio de ingeniería que vale la pena conservar: no todos los problemas de integración deben empaquetarse de inmediato como un nuevo MCP Server. La estandarización del protocolo reduce el costo de integración, pero no resuelve automáticamente los problemas de diseño de permisos, división de tareas y observabilidad.

\begin{warningbox}{MCP no sustituye al diseño de sistemas}
\begin{itemize}
    \item Si la API subyacente ya es inestable, MCP solo expone esa inestabilidad de forma más estandarizada;
    \item Si los límites de la tarea no son claros, aunque el modelo pueda invocar herramientas, la automatización puede terminar en el nivel equivocado;
    \item Si faltan mecanismos de registro, reintentos y reversión, la unificación del protocolo puede incluso amplificar el alcance de propagación de los errores.
\end{itemize}
\end{warningbox}

Esto también explica por qué el curso, mientras explica el protocolo, sigue insistiendo en el agent loop, el context management y las restricciones de producto. Un Coding Agent verdaderamente confiable nunca es «cuantas más herramientas, mejor», sino «capaz de completar la tarea de manera estable dentro de límites claros».
